\chapter{INTRODUCCIÓN DEL CASO}

Amazon.com Inc. es una de las mayores corporaciones tecnológicas del mundo, reconocida por su fuerte cultura de automatización y optimización basada en datos, aplicada tanto a sus operaciones logísticas como a la gestión de recursos humanos. Entre 2014 y 2017, un equipo de especialistas en aprendizaje automático de la compañía, con sede en Edimburgo, desarrolló de manera experimental una herramienta de inteligencia artificial destinada a automatizar la revisión de currículums y la selección de candidatos para puestos técnicos, principalmente de desarrollo de software \cite{dastin2018}.

Los actores implicados en este caso incluyen a Amazon como organización responsable del tratamiento y desarrollo del sistema; los candidatos —mujeres, en particular— cuyas solicitudes de empleo fueron evaluadas y, en muchos casos, penalizadas de forma indirecta por el algoritmo; el periodista Jeffrey Dastin y la agencia Reuters, quienes revelaron públicamente el caso en octubre de 2018; y, de forma indirecta, los organismos reguladores —como la Comisión para la Igualdad de Oportunidades en el Empleo de Estados Unidos (EEOC) y, posteriormente, la Unión Europea— cuyos marcos normativos resultan hoy directamente aplicables a situaciones de esta naturaleza \cite{chen2023}.

La cronología básica del caso puede resumirse así: en 2014, Amazon inició el desarrollo del sistema con el objetivo de mecanizar la búsqueda de talento, otorgando a cada candidato una puntuación de una a cinco estrellas, de forma similar a como se califican los productos en su plataforma de comercio electrónico \cite{dastin2018}. Hacia 2015, el equipo descubrió que el sistema no evaluaba a los candidatos de forma neutral en cuanto al género, penalizando sistemáticamente a las mujeres en los puestos técnicos. A pesar de los intentos de corrección, el proyecto fue finalmente abandonado en 2017, sin llegar a utilizarse de forma exclusiva o generalizada en los procesos reales de contratación \cite{infobae2018}. El caso permaneció confidencial hasta que Reuters lo hizo público en 2018, momento en el que se convirtió en uno de los ejemplos más citados de discriminación algorítmica en el ámbito laboral \cite{chen2023,kochling2020}.
